\documentclass[10pt,a4paper]{amsart}
\usepackage[margin=19mm]{geometry}
\usepackage{amsmath,amssymb,mathtools}
\usepackage[scr=rsfs,cal=euler]{mathalfa}
\usepackage{xfrac}
\usepackage[unicode,hidelinks,bookmarks=false]{hyperref}
\newcommand{\ie}{i.e.\ }
\newcommand{\cf}{cf.\ }
\newcommand{\resp}{respectively\ }
\newcommand{\Subsection}{\subsection}
\providecommand{\nobreakdash}{\nobreak\hbox{-}}
\setlength{\emergencystretch}{2em}
\multlinegap0pt
\usepackage{bm}
\usepackage{bbm}
\usepackage{dsfont}
\usepackage{xspace}
\usepackage{cancel}
\usepackage{mathtools}
\usepackage{amsthm}
\makeatletter
\@ifundefined{lem}{\newtheorem{lem}{Lem}}{}
\makeatother
\title[Quantum K-Rings of Partial Flag Varieties, Coulomb Branches,]{Quantum K-Rings of Partial Flag Varieties, Coulomb Branches, and the Bethe Ansatz}
\author{Irit Huq-Kuruvilla}
\date{Source: arXiv:2409.15575v9 (CC BY 4.0)}
\begin{document}
\maketitle

\noindent\emph{Source and licence.} Quantum K-Rings of Partial Flag Varieties, Coulomb Branches, and the Bethe Ansatz. Version arXiv:2409.15575v9 (CC BY 4.0), distributed under the Creative Commons Attribution 4.0 International licence, as recorded in the arXiv licence metadata for this exact version and shown on the abstract page https://arxiv.org/abs/2409.15575v9. Full rights notice: licenses/arxiv-2409.15575-CC-BY-4.0.txt.\par\smallskip

\noindent\textit{Editorial scope.} Counted unit: the Lem whose source label is \texttt{None} in the article (source0.tex lines 652--662, source label \texttt{None}), delivered together with the article's own complete original proof of it (source0.tex lines 664--763, one \texttt{proof} environment of 100 lines). No mathematics is written, repaired or re-derived: statement and proof are the article's own text line for line. Required context retained uncounted: none. Every definition, display and label of those lines is retained unchanged. Omitted from this package and not redistributed: 1--651, 663--663, 764--1031. Nothing inside a retained excerpt is omitted or altered: every retained body is delivered exactly as the article's lines are written, so no omission receipt of any kind accompanies this package. Citations: the retained excerpts carry no citation command at all, so no bibliography entry of the article is transcribed and no reference is redistributed. Reference adaptation: none was needed and none was made; every reference inside the delivered unit resolves inside it, so no unresolved reference or citation is produced. Printed slips kept literal: no printed word, symbol, hypothesis or number of the counted statement or of its proof was altered. Layout and class: the article's own class is \texttt{amsart} with packages including graphicx, amsmath, hyperref, fullpage, amssymb, amsthm, amsfonts, enumerate, textcomp, eurosym, epsfig, epstopdf, tikz, esint; the wrapper is the lane's pinned \texttt{amsart} editorial wrapper at 10pt on A4, which restates the theorem-like environments the retained text needs and adds the article's own macro definitions that the retained text needs (none), with extra wrapper packages (bm, bbm, dsfont, xspace, cancel, mathtools). The delivered numbering is the unit's own. Assets: no retained span contains a figure environment, a graphics-inclusion command, a PDF-page inclusion, a table or a TikZ picture; no image file is copied into this package, the package declares no asset dependency and no image file is redistributed with it. Licence: version arXiv:2409.15575v9 of this article is distributed under the Creative Commons Attribution 4.0 International licence, as recorded in the arXiv licence metadata for this exact version and shown on the abstract page https://arxiv.org/abs/2409.15575v9. A full rights notice is delivered at licenses/arxiv-2409.15575-CC-BY-4.0.txt. Characters: every retained excerpt is pure ASCII, so the delivered engine, XeLaTeX with its default Latin Modern text font, reports no missing character.
\begin{lem}
For an integer $0\leq \ell \leq v_{i+1}-v_i$, let $F_\ell(d^{i})$ denote the maximum sum of $\ell$ distinct $d^i_j$s.

For any choice of $i,j$, we have: 

$$deg(J_d)< -F_\ell(d^{i+1})+(v_{i}-v_{i+1}+\ell)d^i_j$$
    
    


\end{lem}

\begin{proof}

We first address the case $\ell=0$, corresponding to the inequality:

$$deg(J_d)< (v_{i}-v_{i+1})max_j(d^i_j)$$


As before, we assume without loss of generality that $d^n_j$ are not all equal to 0.

Let $b_{i}=max_j(d^{i}_j)$. 

We have that for all $j$, 

\begin{equation}
    \label{master}
    \sum_j sum_{k=v_{i+1}-v_i}^{v_{i+1}} max(0,d^{i}_j-d^{i+1}_k)\leq R_i. 
\end{equation}

Thus, choosing the $j$ which maximizes the value of $d^i_j$, recalling that $T_i\geq 0$:
\begin{equation}
    \label{grn1}
    T_i-R_i\leq \sum_{k=(v_{i}-v_{i+1})}^{v_{i+1}}min(0,d^{i+1}_k-d^i_j)
\end{equation}


We can further obtain the following bound on the next term in the expression for the degree:

\begin{equation}
    \label{grn2}
    T_{i+1}-R_{i+1}\leq \sum_{k=(v_{i}-v_{i+1})}^{v_{i+1}}d^{i+2}_k-d^{i+1}_k
\end{equation}

This inequality holds because we can extract terms $-\binom{d^{i+1}_k-d^{i+2}_k+1}{2}$ from $T_{i+1}$.

The same inequality holds for $i+s$ for arbitrary $s$. In addition, we note that it is necessarily strict at some point (eventually all the $d^{i+s}_k$ must vanish).

For a given $k$, the contribution to the sum of all of these inequalities is the following:


$$min(0,d^{i+1}_k-d^i_j)+(d^{i+1}_k-d^{i+2}_k)+\dots d^n_k$$



If we eliminate the $min(0,\cdot)$s from the sum, the right hand side telescopes to $(v_i-v_{i+1})d^i_j$, yielding the desired inequality. In fact, the same inequality holds with the $min(0,)$ present. The minimum being achieved at $0$ of a given term corresponds to $d^{s+1}_{k+v_{s+1}-v_{i+1}}\geq d^{s}_{k+v_s-v_{i+1}}$. 

This means that so long as at least one of the terms in the sum for a fixed $k$ is nonzero, which is guaranteed since $d^{n+1}_k=0$ for all $k$, the sum of those terms telescopes to a quantity that must be at most $-d^i_j$, proving the desired inequality. 









%We note that this inequality is necessarily strict for $i=n$, since $b^{n+1}=0$, but by our assumption, $b_n\neq 0$. 

%Moving to the next term. We can use inequality \eqref{master} to write: 

%\begin{equation}
 %   \label{blu}
    %T_{i+1}-R_{i+1}\leq \sum_{k=1}^{v_{i+1}}\sum_{w=v_{i+2}-v_{i+1}}^{v_{i+2}}min(0,d^{i+2}_w-d^{i+1}_k)
%\end{equation}

%If we make the assumption that $v_{i+2}-v_{i+1}>v_{i+1}-v_i$,  we can choose terms from the right side of the inequality, one for each $k$ appearing in inequality \eqref{grn}: 

 %$$T_{i+1}-R_{i+1}\leq \sum_{k=v_{i+1}-v_i}^{v_i}min(0,d^{i+2}_{k+v_{i+2}-v_{i+1}}-d^{i+1}_k)$$

%We can continue to do this procedure for $i$ up to $n$. 

%Taking the sum of these inequalities from $i$ to $n+1$ (noting that the last one is strict), we get:

%$$deg(J_d)<\sum_{k=v_{i+1}-v_i}^{v_{i+1}} min(0,(d^{i+1}_k-d^i_j)+\sum_{s=i+1}^{n} min(0,d^{s+1}_{k+v_{s+1}-v_{i+1}}-d^{s}_{k+v_s-v_{i+1}})$$



%In fact, the assumption that the $v_{i+1}-v_i$s form an increasing sequence is not actually necessary. We can prove the stronger inequality:


 %$$T_{i+1}-R_{i+1}\leq \sum_{k=v_{i+1}-v_i}^{v_i}min(0,d^{i+2}_{k}-d^{i+1}_k)$$

%If $k>v_{i+1}$, the term $min(0,d^{i+2}_{k}-d^{i+1}_k)$ appears as part of $R_{i+1}$. If not, and it is nonzero, we can prove a stronger bound on $T_{i+1}$. If $d^{i+2}_j<d^{i+1}_k$, then we have $$T_i \leq \binom{d^{i+1}_k-d^{i+1}_j+1}{2}-\binom{d^{i+1}_k-d^{i+2}_j+1}{2}\leq d^{i+1}_k-d^{i+2}_j












For general $\ell$. If $F_\ell(d^{i+1})+(v_{i+1}-v_i-\ell)max_j(d^i_j)$ is smaller than the corresponding term for $\ell=0$, the same bound applies. If it is larger, that means the largest $\ell$ $d^{i+1}_k$s total to more than $\ell max_j(d^i_j)$. We can replace $\ell$ of the $d^i_j$s with these $d^{i+1}_k$s and run the same argument.  




\end{proof}

\end{document}
